\documentclass[a4paper,10pt,fleqn]{article}
\usepackage[margin=1in]{geometry}
\usepackage{amssymb}
\usepackage{adjustbox}
\usepackage{natbib}
\usepackage[colorlinks=true,linkcolor=blue,citecolor=blue,urlcolor=blue]{hyperref}
\usepackage{fuzz}
\usepackage{schemapk}
\usepackage{zed-maths}
\usepackage{zed-proof}
\newdimen\savedleftskip
\begin{document}

\section*{Schema Conjunction via Inclusion}
\addcontentsline{toc}{section}{Schema Conjunction via Inclusion}

\noindent Define two simple schemas.

\bigskip

\begin{schema}{A}
x : \nat
\end{schema}

\begin{schema}{B}
y : \nat
\end{schema}

\noindent Schema AB includes both A and B, adds a third component, and uses the schemas as predicates in the where clause.

\bigskip

\begin{schema}{AB}
A \\
B \\
z : \nat
\where
A \land B \\
z = x + y
\end{schema}

\section*{Typed Disambiguation}
\addcontentsline{toc}{section}{Typed Disambiguation}

\noindent A critical disambiguation: count, count' : N is a typed declaration (two variables sharing a type), NOT an inclusion of schema count. The scan-ahead rule detects the colon and routes to the typed-declaration path.

\bigskip

\begin{schema}{TwoNames}
count : \nat \\
limit : \nat
\where
count \leq limit
\end{schema}

\noindent But Counter alone on its own line (no colon) is a bare schema inclusion:

\bigskip

\begin{schema}{HasCounter}
TwoNames \\
extra : \nat
\end{schema}

\end{document}